% ============================================================
%  LAB1 -- Linear Programming Assignment (airline problem)
%  Data Analysis, Optimization and Decision Making -- UAB 2026-27
%  Deadline: Wednesday 11 November 2026, 23:59 (Campus Virtual)
%  Template: IEEE Access (ieeeaccess.cls)
%  Figures: MATLAB scripts in ../matlab/ -> ./figures/
%  Auto-compile: run   latexmk -pvc   in this folder
% ============================================================
\documentclass{ieeeaccess}
\usepackage{amsmath,amssymb,amsfonts}
\usepackage{graphicx}
\graphicspath{{figures/}}
\usepackage{textcomp}
\usepackage{xcolor}
\AtBeginDocument{\renewcommand{\ttdefault}{lmtt}}  % Courier not installed
\newcommand{\todo}[1]{\textcolor{red}{\textbf{TODO:} #1}}

\begin{document}
\history{Date of submission 11 November 2026.}
\doi{--}

\title{Linear Programming Assignment: Optimal Number of Flights for Two Aircraft}
\author{\uppercase{Mats Joris}\authorrefmark{1},
\todo{Author 2 (add name)}\authorrefmark{1}, and
\todo{Author 3 (add name)}\authorrefmark{1}}

\address[1]{Master on Research and Innovation in Computer based Science and
Engineering, Universitat Aut\`onoma de Barcelona, Spain
(NIU: 1826859, \todo{NIU author 2}, \todo{NIU author 3})}

\markboth
{Joris \headeretal: Linear Programming Assignment}
{Joris \headeretal: Linear Programming Assignment}

\corresp{Corresponding author: Mats Joris (NIU: 1826859).}

\begin{abstract}
An airline has two aircraft and wants to know how many flights each one
should make. We model this as a linear programming problem and solve it with
the graphical method and with the simplex method in MATLAB. The maximum
benefit is 190\,000 euro with 160 flights of A and 40 flights of B. The
minimum fuel use is with 30 flights of each aircraft.
\end{abstract}

\begin{keywords}
Graphical method, linear programming, MATLAB, simplex method.
\end{keywords}

\titlepgskip=-21pt

\maketitle

% ============================================================
\section{Introduction}
% ============================================================
An airline has two aircraft, A and B, for one route. We want to know how many
flights each aircraft should make to earn the most money. We write this as a
linear programming (LP) problem and solve it.

% ============================================================
\section{Problem Formulation}
% ============================================================
Let $x_1$ be the number of flights of aircraft A and $x_2$ the number of
flights of aircraft B. The objective function is
\begin{equation}
  \max Z = 1000\,x_1 + 750\,x_2 .
  \label{eq:objective}
\end{equation}
The constraints are:
\begin{align*}
  x_1 &\ge x_2           && \text{(C1)}\\
  x_1 - x_2 &\le 120     && \text{(C2)}\\
  x_2 &\le 180           && \text{(C3)}\\
  x_1 + x_2 &\ge 60      && \text{(C4)}\\
  x_1 + x_2 &\le 200     && \text{(C5)}\\
  x_1,\ x_2 &\ge 0       && \text{(nonnegativity)}
\end{align*}
For C1 we use $x_1 \ge x_2$ instead of $x_1 > x_2$, because LP does not
work with strict inequalities.

% ============================================================
\section{Part 1a: Graphical Solution}
% ============================================================
We draw every constraint as a line (Fig.~\ref{fig:graphical}). The green
area is the feasible region. Then we compute $Z$ in each extreme point
(Table~\ref{tab:points}).

\Figure[t!](topskip=0pt, botskip=0pt, midskip=0pt)[width=0.95\columnwidth]{fig_part1a_graphical.pdf}
{Graphical solution. The dashed line is the isoprofit line through the optimum.
\label{fig:graphical}}

\begin{table}[h]
\centering
\caption{Extreme points and value of $Z$.}
\label{tab:points}
\begin{tabular}{cccc}
\hline
Point & Intersection & $(x_1, x_2)$ & $Z$ (\texteuro) \\
\hline
A & C4 and $x_2 = 0$ & $(60, 0)$    & 60\,000  \\
B & C2 and $x_2 = 0$ & $(120, 0)$   & 120\,000 \\
C & C2 and C5        & $(160, 40)$  & \textbf{190\,000} \\
D & C1 and C5        & $(100, 100)$ & 175\,000 \\
E & C1 and C4        & $(30, 30)$   & 52\,500  \\
\hline
\end{tabular}
\end{table}

$Z$ is highest in point C. So the best solution is $x_1 = 160$ flights for
A and $x_2 = 40$ flights for B, which gives $Z = 190\,000$\,\texteuro.

% ============================================================
\section{Part 1b: Redundant Constraints}
% ============================================================
Yes. Constraint C3 ($x_2 \le 180$) is a redundant constraint, because it
does not change the feasible region. From C1 and C5 we get
$2x_2 \le x_1 + x_2 \le 200$, so $x_2 \le 100$ in the whole feasible region.
In Fig.~\ref{fig:graphical} the line of C3 lies above the green area. If we
remove C3, the feasible region and the solution stay the same.

C1 and C4 are also not used at the optimum (point C is on C2 and C5). If we
remove one or both of them, the feasible region gets bigger, but point C is still the
best point. Only C2 and C5 cannot be removed, because they define point C.

% ============================================================
\section{Part 1c: Simplex Method with MATLAB}
% ============================================================
MATLAB solves LP problems with \texttt{linprog}, in the form
\[
  \min_x f^T x \quad \text{subject to} \quad A x \le b,\ \ x \ge lb .
\]
Our problem is a maximization, so we minimize $-Z$:
\[
  \min\, -Z = -1000\,x_1 - 750\,x_2 .
\]
All constraints are written as $\le$. For C1 and C4 we multiply by $-1$:
\[
  A = \begin{bmatrix} -1 & 1 \\ 1 & -1 \\ 0 & 1 \\ -1 & -1 \\ 1 & 1 \end{bmatrix},\quad
  b = \begin{bmatrix} 0 \\ 120 \\ 180 \\ -60 \\ 200 \end{bmatrix},
\]
\[
  f = \begin{bmatrix} -1000 \\ -750 \end{bmatrix},\quad
  lb = \begin{bmatrix} 0 \\ 0 \end{bmatrix}.
\]
We use the simplex algorithm of \texttt{linprog}:
\begin{verbatim}
options = optimoptions('linprog', ...
          'Algorithm','dual-simplex');
[x,fval] = linprog(f,A,b,[],[],lb,[],options)
Z = -fval
\end{verbatim}
MATLAB gives:
\begin{verbatim}
Optimal solution found.
x =
   160
    40
fval =
     -190000
Z =
      190000
\end{verbatim}
So $x_1 = 160$, $x_2 = 40$ and $Z = -\text{fval} = 190\,000$\,\texteuro.
This is the same result as with the graphical method (point C).

% ============================================================
\section{Part 2: Minimum Fuel Consumption}
% ============================================================
Let one flight of B use 1 unit of fuel. Then one flight of A uses 2 units.
The constraints are the same as in Part 1, but now we minimize the total fuel:
\[
  \min F = 2\,x_1 + x_2 .
\]

\textit{Graphical method.} The feasible region is the same as in Part 1
(Fig.~\ref{fig:fuel}). We compute $F$ in each extreme point
(Table~\ref{tab:fuel}).

\Figure[t!](topskip=0pt, botskip=0pt, midskip=0pt)[width=0.95\columnwidth]{fig_part2_fuel.pdf}
{Graphical solution of Part 2. The dashed line is the isocost line through
the optimum.\label{fig:fuel}}

\begin{table}[h]
\centering
\caption{Extreme points and value of $F$.}
\label{tab:fuel}
\begin{tabular}{ccc}
\hline
Point & $(x_1, x_2)$ & $F$ (units) \\
\hline
A & $(60, 0)$    & 120 \\
B & $(120, 0)$   & 240 \\
C & $(160, 40)$  & 360 \\
D & $(100, 100)$ & 300 \\
E & $(30, 30)$   & \textbf{90} \\
\hline
\end{tabular}
\end{table}

$F$ is lowest in point E, so $x_1 = 30$, $x_2 = 30$ and $F = 90$ units.

\textit{Simplex method with MATLAB.} We use the same $A$, $b$ and $lb$ as in
Part 1c, with $f = [2 \;\; 1]^T$. MATLAB gives:
\begin{verbatim}
Optimal solution found.
x =
    30
    30
fval =
    90
\end{verbatim}
This is the same result as with the graphical method.

\textit{Note.} In point E both aircraft make 30 flights, so A does not fly
more than B. If A must make at least one flight more ($x_1 \ge x_2 + 1$), the
best whole numbers are $x_1 = 31$ and $x_2 = 29$, with $F = 91$ units.

% ============================================================
\section{Conclusion}
% ============================================================
The maximum benefit is $Z = 190\,000$\,\texteuro{} with $x_1 = 160$ and
$x_2 = 40$. The graphical method and the simplex method in MATLAB give the
same result. Constraint C3 is redundant. The minimum fuel use is $F = 90$
units with $x_1 = x_2 = 30$ (or $F = 91$ with $x_1 = 31$, $x_2 = 29$ if A
must fly strictly more than B).

\EOD
\end{document}
